% ============================================================
%  BAB 6 – PENUTUP
%  Compile standalone:  pdflatex bab-6.tex
%  Compile as part of thesis: pdflatex main.tex
% ============================================================
\documentclass[main]{subfiles}
\usepackage{hyphenat}
\usepackage{iftex}

% ── Encoding, Language & Fonts ───────────────────────────────
\ifXeTeX
    \usepackage{polyglossia}
    \setmainlanguage{bahasai}
    \usepackage{fontspec}
    \setmainfont{TeX Gyre Termes}
    \setsansfont{TeX Gyre Heros}
    \setmonofont[Scale=0.88]{TeX Gyre Cursor}
\else
    \usepackage[utf8]{inputenc}
    \usepackage[T1]{fontenc}
    \usepackage[indonesian]{babel}
    \usepackage{mathptmx}
    \usepackage{helvet}
    \usepackage{courier}
\fi

% ── Page geometry (A4, UI thesis margins) ────────────────────
\usepackage[
  a4paper,
  top    = 2.54cm,
  bottom = 2.54cm,
  left   = 3.17cm,
  right  = 2.54cm
]{geometry}

% ── Line spacing & paragraph ─────────────────────────────────
\usepackage{setspace}
\onehalfspacing
\setlength{\parindent}{1.27cm}
\setlength{\parskip}{0pt}

% ── Microtype & justification ────────────────────────────────
\usepackage{microtype}
\sloppy

% ── Headings ─────────────────────────────────────────────────
\usepackage{titlesec}

\titleformat{\chapter}[block]
  {\centering\sffamily\bfseries\large}{}
  {0pt}{\MakeUppercase}
\titlespacing*{\chapter}{0pt}{0pt}{1.5em}

\titleformat{\section}
  {\sffamily\bfseries\normalsize}{\thesection}{1em}{}
\titlespacing*{\section}{0pt}{1.2em}{0.6em}

\titleformat{\subsection}
  {\sffamily\bfseries\normalsize}{\thesubsection}{1em}{}
\titlespacing*{\subsection}{0pt}{1em}{0.5em}

\titleformat{\subsubsection}
  {\sffamily\bfseries\normalsize}{\thesubsubsection}{1em}{}
\titlespacing*{\subsubsection}{0pt}{0.8em}{0.4em}

% ── Math ─────────────────────────────────────────────────────
\usepackage{amsmath, amssymb}

% ── Lists ────────────────────────────────────────────────────
\usepackage{enumitem}
\setlist[enumerate,1]{
  label      = \arabic*.,
  leftmargin = 1.27cm,
  itemsep    = 4pt,
  parsep     = 0pt,
  topsep     = 4pt
}
\setlist[enumerate,2]{
  label      = \alph*.,
  leftmargin = 1.27cm,
  itemsep    = 3pt,
  parsep     = 0pt,
  topsep     = 2pt
}

% ── Hyperlinks ───────────────────────────────────────────────
\usepackage[hidelinks]{hyperref}

% Bibliography setup is inherited from main.tex (biblatex + references.bib)
% when this file is compiled standalone via the subfiles class.

% ============================================================
\begin{document}

% When standalone: set counter so \chapter{} renders as "BAB 6".
% When in main.tex: main.tex sets the counter before \subfile{bab-6}.
\ifSubfilesClassLoaded{}{\setcounter{chapter}{5}}

\chapter{Penutup}\label{sec:penutup}

Bab ini menutup laporan penelitian dengan menarik kesimpulan atas keseluruhan
proses perancangan, pengembangan, dan evaluasi artefak yang telah dipaparkan
pada bab-bab sebelumnya. Kesimpulan disusun untuk menjawab secara langsung
keempat pertanyaan penelitian yang dirumuskan pada Subbab~\ref{sec:ruang-lingkup},
dilanjutkan dengan pemaparan keterbatasan penelitian dan saran bagi
pengembangan lanjutan.

% ============================================================
\section{Kesimpulan}
% ============================================================

Penelitian ini, dalam kerangka \textit{Design Science Research} (DSR), berhasil
merancang, membangun, dan mengevaluasi sebuah \textit{pipeline}
\textit{LLM-Assisted Knowledge Graph Completion} (KGC) untuk memetakan
interkoneksi materi sains lintas-disiplin (Fisika, Kimia, dan Biologi) pada
Fase~F Kelas~XII Kurikulum Merdeka berbahasa Indonesia. Berdasarkan hasil
evaluasi struktural dan validasi pakar yang dibahas pada Bab~5, kesimpulan yang
menjawab keempat rumusan masalah penelitian adalah sebagai berikut.

\begin{enumerate}
  \item \textbf{Pemetaan interkoneksi untuk mengatasi \textit{siloed learning}.}
        \textit{Pipeline} yang dikembangkan mampu memetakan interkoneksi semantik
        antar-materi sains secara eksplisit, sehingga menjawab masalah
        fragmentasi kurikulum. Sebelum KGC, ketiga mata pelajaran membentuk
        komunitas yang nyaris terisolasi, ditandai nilai \textit{Modularity}
        0{,}811 yang merefleksikan kondisi \textit{siloed learning}. Setelah
        penambahan 125 relasi lintas-buku, \textit{Modularity} turun menjadi
        0{,}373, sementara \textit{Average Degree Centrality} (ADC) meningkat
        98{,}4\% (dari 0{,}74 menjadi 1{,}47) dan \textit{Graph Density}
        antar-konsep hampir berlipat ganda (dari 0{,}00108 menjadi 0{,}00214).
        Bukti topologis ini (Tabel~\ref{tab:struktural}) mengonfirmasi bahwa
        \textit{Knowledge Graph} (KG) yang dihasilkan berhasil menjembatani
        ketiga komunitas disiplin yang semula terpisah.

  \item \textbf{Penanganan keterbatasan metode konvensional pada entitas
        \textit{long-tail}.} Alih-alih mengandalkan KGC berbasis \textit{embedding}
        struktural yang terbukti menurun performanya pada entitas \textit{long-tail},
        penelitian ini menempuh pendekatan semantik dua tahap
        (Subbab~\ref{sec:graph-completion}): pencarian kandidat melalui ANN
        \textit{vector index} (HNSW) atas \textit{embedding} deskripsi konsep
        (\texttt{gemini-embedding-001}), dilanjutkan klasifikasi relasi oleh LLM
        dengan strategi \textit{few-shot In-Context Learning}. Karena pembangkitan
        kandidat bertumpu pada kemiripan makna deskripsi---bukan pada pola
        konektivitas graf---konsep yang berkoneksi sedikit (\textit{long-tail})
        tetap memperoleh kandidat relasi yang bermakna. Pendekatan ini terbukti
        mampu menemukan relasi lintas-disiplin yang valid (validitas 99{,}1\%)
        yang tidak akan tertangkap oleh metode struktural murni.

  \item \textbf{Efektivitas LLM-Assisted KGC untuk relasi implisit
        lintas-disiplin.} LLM-Assisted KGC sangat efektif dalam \textit{menemukan}
        dan \textit{melabeli} relasi implisit, namun memiliki kelemahan sistematis
        pada penentuan arah relasi. Pada validasi pakar Fase~2 atas relasi
        lintas-buku, sistem mencapai validitas 99{,}1\% (116/117) dan ketepatan
        tipe 98{,}3\% (115/117), tetapi ketepatan arah hanya 75{,}9\% (88/116).
        Kesalahan arah bersifat sistematis: 26 dari 28 relasi terbalik berasal
        dari pasangan Biologi\,$\leftrightarrow$\,Kimia, dengan LLM cenderung
        keliru menempatkan konsep Biologi sebagai sumber. Secara semantik, jenis
        relasi didominasi \texttt{PRASYARAT\_UNTUK} (45{,}6\%), menandakan bahwa
        jembatan lintas-disiplin pada kurikulum kelas XII lebih banyak bersifat
        prasyarat pengetahuan, dengan Kimia berperan sebagai poros (\textit{hub})
        yang menghubungkan Biologi (78\% relasi) dan Fisika
        (Gambar~\ref{fig:bridge-graph}).

  \item \textbf{Keandalan evaluasi hybrid dan pengukuran kualitas kuantitatif.}
        \textit{Pipeline} evaluasi hybrid berhasil mengukur kualitas KG secara
        andal dan terkuantifikasi. Pada validasi intra-mapel (Fase~1) oleh enam
        pakar domain, sistem mencatatkan \textit{precision} 93{,}6\%,
        \textit{recall} 91{,}7\%, dan F1-Score 92{,}6\%
        (Tabel~\ref{tab:validasi-metrik}), dengan kesepakatan antarpakar Gwet's
        AC1 yang bervariasi dari ``baik'' hingga ``sangat baik'' (0{,}607 pada
        Kimia hingga 0{,}931 pada Biologi). Mekanisme \textit{LLM-as-a-judge}
        (Gemini~3.1 Pro Preview) berfungsi sebagai penengah yang menyelesaikan
        ketidaksepakatan intra-mapel berbasis konteks buku teks, sementara
        validasi panel pada relasi lintas-disiplin (Fase~2) menghasilkan
        validitas 99{,}1\%. Dengan demikian, kombinasi metrik struktural dan
        validasi pakar berbantuan LLM terbukti mampu mempertanggungjawabkan
        validitas keilmuan KG secara objektif dan terukur.
\end{enumerate}

Secara keseluruhan, penelitian ini membuktikan bahwa pendekatan LLM-Assisted KGC
merupakan metode yang layak (\textit{viable}) untuk membangun dan melengkapi graf
pengetahuan kurikulum sains lintas-disiplin berbahasa Indonesia. Artefak yang
dihasilkan---berupa purwarupa KG tervalidasi beserta \textit{pipeline}-nya yang
terdokumentasi---memenuhi tujuan penelitian dan menyediakan fondasi data terbuka
pertama bagi pengembang kurikulum dan peneliti teknologi pendidikan, dengan catatan
penting bahwa kanonikalisasi arah relasi masih memerlukan penyempurnaan lebih lanjut.

% ============================================================
\section{Keterbatasan Penelitian}
% ============================================================

Penelitian ini memiliki sejumlah keterbatasan yang perlu diperhatikan dalam
menafsirkan hasilnya.

\begin{enumerate}
  \item \textbf{Cakupan korpus.} Korpus terbatas pada buku teks siswa Fisika,
        Kimia, dan Biologi Fase~F Kelas~XII (\textit{e-book} resmi
        Kemendikbudristek). Buku teks Kelas~XI dan dokumen Capaian Pembelajaran
        (CP) tidak disertakan, sehingga generalisasi temuan ke fase atau jenjang
        lain belum dapat dipastikan.

  \item \textbf{Kesalahan arah relasi.} Sebagaimana disimpulkan di atas, sekitar
        seperempat relasi lintas-buku memiliki arah yang terbalik. Analisis
        \textit{parameter sweep} (Tabel~\ref{tab:validasi-sweep}) menunjukkan bahwa
        kesalahan ini tidak teratasi dengan pengetatan \textit{threshold},
        melainkan memerlukan perbaikan pada \textit{prompt} klasifikasi atau
        kanonikalisasi arah pasca-proses yang belum diimplementasikan pada
        penelitian ini.

  \item \textbf{\textit{Under-extraction} relasi kuantitatif.} LLM cenderung kuat
        pada konsep deskriptif namun lemah pada relasi berbasis persamaan
        matematis; 67\% \textit{missing triples} yang diusulkan pakar Fisika
        berupa rumus yang tidak berhasil diekstrak
        (Subbab~\ref{subsubsec:error-analysis}).

  \item \textbf{Batasan inferensi pedagogis.} Sistem hanya menginferensi relasi
        \textit{single-hop} dan tidak menyusun jalur belajar multi-langkah, tidak
        memodelkan tingkat kognitif (taksonomi Bloom) maupun keadaan pembelajar
        (\textit{learner state}).

  \item \textbf{Keterbatasan metrik dan model.} \textit{Recall} dan F1 pada
        validasi lintas-disiplin (Fase~2) tidak dapat dihitung karena tiadanya
        \textit{denominator} relasi benar yang menyeluruh dalam konfigurasi panel
        tunggal. Selain itu, penelitian hanya memanfaatkan satu keluarga model
        (Gemini) tanpa komparasi lintas-model, dan graf diinstansiasi pada skala
        purwarupa (Neo4j Aura \textit{Free tier}) tanpa integrasi ke sistem
        pengguna akhir.
\end{enumerate}

% ============================================================
\section{Saran}
% ============================================================

Berdasarkan kesimpulan dan keterbatasan di atas, beberapa saran diajukan untuk
penelitian dan pengembangan selanjutnya.

\subsection{Bagi Peneliti dan Pengembang Teknologi Pendidikan}

\begin{enumerate}
  \item \textbf{Penanganan arah relasi.} Memprioritaskan perbaikan ketepatan arah,
        misalnya melalui penambahan contoh \textit{few-shot} yang eksplisit
        menunjukkan arah, langkah klasifikasi arah terpisah, atau kanonikalisasi
        pasca-proses (misalnya memberlakukan kecenderungan arah Kimia$\to$Biologi
        sebagai prasyarat) untuk mengoreksi 28 relasi terbalik secara sistematis.

  \item \textbf{Penguatan ekstraksi relasi kuantitatif.} Menambahkan tahap
        ekstraksi khusus untuk rumus dan persamaan matematis, atau memanfaatkan
        kemampuan \textit{multimodal} model untuk membaca notasi matematis dan
        diagram, guna mengurangi \textit{under-extraction}.

  \item \textbf{Perluasan korpus dan generalisasi.} Memperluas korpus ke buku teks
        Kelas~XI, dokumen Capaian Pembelajaran (CP), serta mata pelajaran dan fase
        lain untuk menguji generalisasi \textit{pipeline}.

  \item \textbf{Inferensi multi-langkah.} Mengembangkan kemampuan penalaran
        \textit{multi-hop} untuk menyusun jalur belajar lintas-disiplin
        terkomposisi, serta mengintegrasikan taksonomi Bloom dan pemodelan
        \textit{learner state} untuk pengurutan konsep yang adaptif.

  \item \textbf{Evaluasi lintas-model.} Membandingkan beberapa keluarga LLM (mis.\
        GPT, Claude) pada tahap ekstraksi, \textit{completion}, dan
        \textit{LLM-as-a-judge} untuk menguji ketangguhan hasil, serta memperkuat
        validasi Fase~2 dengan dua pakar independen agar \textit{recall} dan AC1
        dapat dihitung.
\end{enumerate}

\subsection{Bagi Pendidik dan Pengembang Kurikulum}

KG yang dihasilkan dapat dimanfaatkan sebagai referensi dalam merancang unit
pembelajaran terpadu lintas-disiplin, khususnya dengan memanfaatkan posisi Kimia
sebagai poros konseptual yang menjembatani Biologi dan Fisika. Pengembangan
lanjutan dapat membangun lapisan aplikasi bagi pengguna akhir di atas KG ini,
seperti perkakas perancangan kurikulum, sistem rekomendasi pembelajaran adaptif,
atau integrasi ke dalam platform LMS, yang berada di luar lingkup penelitian ini.

\ifSubfilesClassLoaded{\printbibliography[title={Daftar Pustaka}]}{}

\end{document}
